\section{Convergence Diagnostics}
\label{sec:diagnostics}

The \re{} chain runner computes three convergence diagnostics in the
same pass as the sampling, at zero additional per-step cost beyond
tracking per-chain statistics.
The diagnostic framework is described in full in~\citet{dellaLibera2026g4};
we summarize the implementations here for completeness.

\subsection{R-hat (Potential Scale Reduction Factor)}

The Gelman-Rubin $\rhat$ statistic~\citep{gelman1992inference} compares
between-chain variance to within-chain variance for a scalar summary
statistic $\theta$ (e.g., Democratic seat count or Polsby-Popper
compactness).
For $C$ chains each of length $N$, let $\bar{\theta}_{c\cdot}$ denote
the within-chain mean and $\bar{\theta}_{\cdot\cdot}$ the grand mean.
The between-chain variance is
\[
  B = \frac{N}{C-1}\sum_{c=1}^C (\bar{\theta}_{c\cdot} - \bar{\theta}_{\cdot\cdot})^2,
\]
and the within-chain variance is
\[
  W = \frac{1}{C}\sum_{c=1}^C s_c^2,
  \quad s_c^2 = \frac{1}{N-1}\sum_{t=1}^N (\theta_{ct} - \bar{\theta}_{c\cdot})^2.
\]
The $\rhat$ statistic is
\[
  \rhat = \sqrt{\frac{(N-1)/N \cdot W + B/N}{W}}.
\]
Values of $\rhat < 1.1$ indicate approximate convergence; we target
$\rhat < 1.05$ for publication-quality results.

For redistricting statistics, which are discrete (seat counts take
integer values) or bounded (compactness metrics lie in $[0,1]$),
we follow \citet{vehtari2021rank} and apply the rank-normalised $\rhat$
variant.
All $CN$ draws across all chains are first rank-transformed to
approximate normality, then $\rhat$ is computed on the transformed values.
Rank normalisation removes sensitivity to the marginal distribution shape
and is the current best practice for MCMC diagnostics~\citep{vehtari2021rank}.

\paragraph{Implementation.}
Each chain accumulates $\bar{\theta}_{c\cdot}$ and $s_c^2$ incrementally
using Welford's online algorithm~\citep{welford1962note}, requiring $O(1)$
storage per chain.
For rank normalisation, all step-level draws are buffered in a
\texttt{Vec<f64>} and rank-sorted at chain completion.
$\rhat$ is computed in a single pass over the $C$ chain statistics.

\subsection{Effective Sample Size}

The effective sample size accounts for autocorrelation within each chain.
Using Geyer's initial monotone sequence estimator~\citep{geyer1992practical},
the ESS for a single chain of length $N$ is
\[
  \ess = \frac{N}{1 + 2\sum_{k=1}^{\hat{K}} \hat{\rho}_k},
\]
where $\hat{\rho}_k$ is the lag-$k$ autocorrelation estimate and
$\hat{K}$ is the cutoff where the monotone sequence estimate becomes
non-positive.
For redistricting chains, typical lag-1 autocorrelation is
$\hat{\rho}_1 \approx 0.85$--$0.95$, giving
$\ess \approx 500$--$2{,}000$ for a 10{,}000-step chain.

We use the multi-chain $\ess$ formula from~\citet{vehtari2021rank}, which
pools within-chain and between-chain information:
\[
  \ess_{\text{bulk}} = \frac{CN \cdot \min(\hat{V}/B, 1)}{1 + 2\sum_{k=1}^{\hat{K}}\hat{\rho}_k^+},
\]
where $\hat{V} = W + B/N$ and $\hat{\rho}_k^+$ is the rank-normalised
autocorrelation at lag $k$.

\paragraph{Target.}
We require $\ess \geq 400$ for any summary statistic used in a legal
filing.
For a 10{,}000-step, 8-chain run at typical redistricting autocorrelation,
the expected $\ess$ is 4{,}000--16{,}000---well above this threshold.

\subsection{Hamming Autocorrelation}

The Hamming distance between two redistricting plans $\sigma$ and $\sigma'$
is the fraction of census tracts assigned to different districts:
\[
  d_H(\sigma, \sigma') = \frac{|\{v : \sigma(v) \neq \sigma'(v)\}|}{n}.
\]
The lag-$k$ Hamming autocorrelation is defined as
\[
  \ham(k) = \mathrm{Corr}(d_H(\sigma_t, \sigma_0), d_H(\sigma_{t+k}, \sigma_0)),
\]
where the correlation is taken across steps $t$.
In practice, $\hat{\rho}_1 \approx 0.87$--$0.93$ for \recom{}, with
geometric decay: $\ham(k) \approx \hat{\rho}_1^k$.

The Hamming autocorrelation serves as a distribution-free mixing
diagnostic: it measures how rapidly consecutive plans decorrelate in
plan space, independent of any particular summary statistic.
A chain with $\ham(1) > 0.98$ is mixing slowly and may require
thinning or a longer run; $\ham(1) < 0.90$ indicates healthy mixing.

The normalized cut fraction $\phi(\sigma) = |E_\sigma| / |E|$
(the fraction of edges crossing district boundaries) is used as a
computationally efficient proxy for plan identity when computing
$\ham(1)$.
This proxy is exact in the sense that $\phi$ is the same summary
statistic minimized by METIS and reported by the \ar{} evaluation.

\subsection{Convergence Targets for Publication}

Table~\ref{tab:diag-targets} summarizes the diagnostic thresholds.
The target is to achieve all three thresholds after 5{,}000 steps for
all six G-track states.
At \re{}'s estimated throughput of 50{,}000 steps/sec, 5{,}000 steps
require approximately 0.1 seconds.

\begin{table}[h]
\centering
\caption{Convergence diagnostic thresholds for \re{} publication results.}
\label{tab:diag-targets}
\begin{tabular}{lll}
\toprule
Diagnostic & Symbol & Target \\
\midrule
Gelman-Rubin (rank-normalised) & $\rhat$ & $< 1.05$ \\
Effective sample size & $\ess$ & $\geq 400$ per statistic \\
Hamming autocorrelation & $\ham(1)$ & $0.85$--$0.95$ \\
\bottomrule
\end{tabular}
\end{table}
